\begin{afterword}
\summaryhead

The author has struggled all their life to keep working over a span of hours, and fills the times when work will not come with web browsing and short videos. Asked by friends why not watch a film instead, the author explained that the ability to work might return at any moment and could not be predicted.

Saying this, the author realized that the friends seemed to assume that other people can predict when they will be able to work again. Perhaps the work cycle was not unusually chaotic; perhaps the author was unusually bad at predicting it. That would be the Mind Projection Fallacy in everyday life: one's own ignorance, seen as chaos in the world. We fail to consider this, the author suspects, not from vanity but because we do not think of ourselves at all. Naming the ignorance does not cure it, but it is a different path from calling the problem chaotic.

\responsehead

This is a short personal note, and it claims little. Nearly every step is hedged (``seemed to be implying'', ``Maybe'', ``I suspect'', ``may or may not help''), and the post says plainly that the new description has not fixed anything. Its use of the Mind Projection Fallacy fits the sense Jaynes gave the term: a state of partial knowledge treated as a property of the thing.

The weak points are small ones. The turn of the post rests on an inference from the friends' words, and as the author reports them those words were conditional: ``if you know you can't work for a while''. That does not say that anyone does know. The author marks the inference with ``seemed'', and the dialogue with ``a very loose gist'', so the post does not lean on it hard, but it is the only evidence that other people can predict their working hours.

The one general claim, that people fail to consider their own ignorance because ``we don't think of ourselves at all'', rests only on the author's case, and is offered with ``I suspect''.

The phrase ``inverted stupidity'' is used as if the reader knew it, with a reference to its use for ``high abstract things like Artificial Intelligence''. It appears nowhere else in the book.

In short: a candid, hedged note that applies the Mind Projection Fallacy to the author's own working habits, with its one general claim about people offered as a suspicion from a single case.
\end{afterword}
